\section{Results with the reference pipeline}
\label{sec:results}

\subsection{Expansion recovers quiet members, and a two-line rule matches it}
\label{sec:results-expansion}

The detector reaches primary-window edge F1 $0.495 \pm 0.027$ and ROC-AUC 0.965.
% source: results/tables/r1_edge_metrics.csv (label=causal); docs/FINDINGS.md (A4)
Table~\ref{tab:ring-methods} gives the ring-level picture; Section~\ref{sec:results-baselines}
discusses the baselines in its upper block.

% Source: results/tables/r9_ring_ci.csv (rows b0, b1, b1b, proposed, average; window=primary) and
%         results/tables/v45_posthoc_ring_ci.csv (baselines and ring-folded rows; produced by
%         evaluation/r9_style_ci.py, which runs evaluation/step9_tables.py unmodified, so every
%         interval in this table is the same ring bootstrap: ring_stats.bootstrap_over_rings).
% Quiet-member recall is the point estimate from r9_quiet_members.csv / v45_posthoc_quiet_members.csv
% (FINDINGS convention, 2026-11-12), not the bootstrap mean in r9_ring_ci.csv.
\begin{table}[t]
\centering
\caption{Ring recovery on the primary window (174 truth rings, 5 seeds or runs,
in-window members, Jaccard $\ge 0.5$). Brackets: 95\,\% bootstrap intervals over truth rings,
seeds averaged inside each resample. Rows marked $^{*}$ use expansion-policy inputs folded by
ring (leak-free); the rows directly above them are the same methods with day-folded inputs, as
first reported. Rows marked $^{\S}$ are post hoc.}
\label{tab:ring-methods}
\small
\begin{tabular}{lccccc}
\toprule
Method & Ring prec. & Ring rec. & Ring F1 [95\,\% CI] & Member cov. & Quiet-member rec. \\
\midrule
CSGM~\citep{tian2025}$^{\dagger\S}$ & 0.321 & 0.048 & \ci{0.083}{0.047}{0.123} & 0.898 & --- \\ % source: v45_posthoc_ring_ci.csv (csgm, primary: 0.3208 / 0.0483 / 0.0834 [0.0474, 0.1227] / 0.8982)
Louvain, resolution 0.5$^{\S}$ & 0.795 & 0.382 & \ci{0.515}{0.428}{0.600} & 0.853 & 0.001 \\ % source: v45_posthoc_ring_ci.csv (louvain_r0.5); quiet v45_posthoc_quiet_members.csv (0.0013)
Louvain, resolution 1$^{\S}$ & 0.748 & 0.358 & \ci{0.483}{0.397}{0.563} & 0.834 & 0.009 \\ % source: v45_posthoc_ring_ci.csv (louvain_r1); quiet 0.0085
Louvain, resolution 2$^{\S}$ & 0.419 & 0.062 & \ci{0.107}{0.070}{0.148} & 0.655 & 0.002 \\ % source: v45_posthoc_ring_ci.csv (louvain_r2); quiet 0.0021
Label propagation$^{\S}$ & 0.785 & 0.377 & \ci{0.509}{0.421}{0.593} & 0.847 & 0.001 \\ % source: v45_posthoc_ring_ci.csv (lpa); quiet 0.0013
\midrule
Seeds + connected components & 0.814 & 0.384 & \ci{0.521}{0.429}{0.607} & 0.866 & 0.005 \\ % source: r9_ring_ci.csv (b0,primary)
Fixed expansion rule    & 0.171 & 0.162 & \ci{0.162}{0.118}{0.213} & 0.970 & 0.118 \\ % source: r9_ring_ci.csv (b1,primary); quiet: r9_quiet_members.csv (b1,primary,ALL = 0.1183)
Rule + score quantile  & 0.702 & 0.463 & \ci{0.555}{0.470}{0.632} & 0.943 & 0.223 \\ % source: r9_ring_ci.csv (b1b,primary); quiet: r9_quiet_members.csv (b1b,primary,ALL = 0.2229)
Learned expansion & 0.838 & 0.446 & \ci{0.582}{0.490}{0.668} & 0.942 & 0.209 \\ % source: r9_ring_ci.csv (proposed,primary); quiet: r9_quiet_members.csv (proposed,primary,ALL = 0.2092)
\quad ring-folded inputs$^{*\S}$ & 0.784 & 0.468 & \ci{0.585}{0.498}{0.670} & 0.960 & 0.251 \\ % source: v45_posthoc_ring_ci.csv (proposed_ringfold: 0.7835 / 0.4678 / 0.5848 [0.4980, 0.6699] / 0.9603); quiet 0.2510
\quad community step instead of components$^{\S}$ & 0.781 & 0.431 & \ci{0.555}{0.466}{0.641} & 0.897 & --- \\ % source: v45_posthoc_ring_ci.csv (lpa_on_expansion: 0.7806 / 0.4310 / 0.5550 [0.4663, 0.6408] / 0.8974)
\textbf{Full (expansion + partition)} & \textbf{0.818} & \textbf{0.461} & \textbf{\ci{0.589}{0.502}{0.675}} & \textbf{0.944} & \textbf{0.203} \\ % source: r4_partition.csv (average) / r9_ring_ci.csv (average,primary)
\quad ring-folded inputs$^{*\S}$ & 0.760 & 0.483 & \ci{0.589}{0.503}{0.672} & 0.960 & 0.243 \\ % source: v45_posthoc_ring_ci.csv (full_ringfold: 0.7597 / 0.4828 / 0.5893 [0.5034, 0.6722] / 0.9604); quiet 0.2425
\bottomrule
\multicolumn{6}{l}{\footnotesize $^{\dagger}$ External baseline, no detector; five runs over random two-institution splits of the accounts.}\\
\multicolumn{6}{l}{\footnotesize Community rows: the method run on the graph of the top-$k$\,\% flagged edges, $k$ chosen on validation per seed.}\\
\multicolumn{6}{l}{\footnotesize ``Community step instead of components'': label propagation on the learned expansion's subgraph, nothing tuned.}\\
\end{tabular}
\end{table}


Expansion from seeds recovers ring members the detector never fires on. With leak-free policy
inputs (Section~\ref{sec:limitations}), quiet-member recall rises from \ci{0.0052}{0.000}{0.015}
with seeds alone to \ci{0.251}{0.179}{0.344} with the learned policy, significant both on quiet
accounts per ring ($+0.432$ [$+0.301$, $+0.590$]) and on whole-ring matching ($+0.084$ [$+0.053$,
$+0.117$]), both $p_{\mathrm{Holm}}=0.0004$.
% source: results/notes/p2_ringfold.md (post hoc table); results/tables/v39_ringfold_quiet_members.csv; v39_ringfold_significance.csv (proposed_ringfold_vs_b0)
With the day-folded inputs we first reported and pre-registered, the recall was
\ci{0.2092}{0.142}{0.290}, with paired differences of $+0.359$ [$+0.236$, $+0.510$] and $+0.062$
[$+0.036$, $+0.092$], both $p_{\mathrm{Holm}}=0.0005$.
% source: results/tables/r9_quiet_members.csv (b0/proposed, primary, ALL); CIs r9_ring_ci.csv; r9_significance.csv (proposed_vs_b0)
It is the only one of the five pre-registered matched-ring comparisons to survive correction
(Table~\ref{tab:significance}), and the result is stable across seed budgets: expansion
recovers 14--24\,\% of quiet members at every budget we tried
(Table~\ref{tab:robustness}, appendix).
% source: docs/FINDINGS.md (section E)

% Source: results/tables/r9_significance.csv (all five pre-registered comparisons)
\begin{table}[t]
\centering
\caption{The five pre-registered comparisons: paired sign-flip permutation tests on per-ring
matched indicators (10{,}000 permutations), Holm-corrected. $\Delta$ is the per-ring matched
rate; the quiet-member column is the per-ring quiet-accounts delta where the family
pre-registered it. The test is blind to precision (Section~\ref{sec:data-stats}).}
\label{tab:significance}
\small
\begin{tabular}{llccc}
\toprule
Comparison & Metric & $\Delta$ [95\,\% CI] & $p_{\mathrm{Holm}}$ & Reject at 0.05 \\
\midrule
proposed vs.\ B0  & matched rate & $+0.062$ [$+0.036$, $+0.092$] & $0.0005$ & yes \\ % source: r9_significance.csv (proposed_vs_b0, mean_delta, p_holm)
proposed vs.\ B0  & quiet accounts/ring & $+0.359$ [$+0.236$, $+0.510$] & $0.0005$ & yes \\ % source: r9_significance.csv (proposed_vs_b0, quiet_mean_delta, quiet_p_holm)
proposed vs.\ B1b & matched rate & $-0.017$ [$-0.041$, $+0.005$] & $0.49$ & no \\ % source: r9_significance.csv (proposed_vs_b1b)
proposed vs.\ B1b & quiet accounts/ring & $-0.024$ [$-0.084$, $+0.037$] & $1.00$ & no \\ % source: r9_significance.csv (proposed_vs_b1b, quiet columns)
full vs.\ expansion only & matched rate & $+0.015$ [$-0.002$, $+0.036$] & $0.49$ & no \\ % source: r9_significance.csv (full_vs_b2)
cc+partition vs.\ B0 & matched rate & $+0.018$ [$+0.002$, $+0.039$] & $0.23$ & no \\ % source: r9_significance.csv (cc_partition_vs_b0)
full vs.\ B1b & matched rate & $-0.002$ [$-0.032$, $+0.028$] & $0.94$ & no \\ % source: r9_significance.csv (full_vs_b1b)
\bottomrule
\end{tabular}
\end{table}


\textbf{The learned policy is not what produces this recovery.} The rule B1b reaches
\ci{0.2229}{0.163}{0.299} quiet recall, indistinguishable from the policy on quiet accounts
($-0.024$ [$-0.084$, $+0.037$], $p_{\mathrm{Holm}}=1.00$) and on whole-ring matching
($p_{\mathrm{Holm}}=0.49$).
% source: docs/FINDINGS.md (B1, B2); r9_significance.csv (proposed_vs_b1b)
We ran this comparison on purpose, against our own method.

What the policy buys is the purity of what it adds. With leak-free policy inputs (cross-scores
folded by ring; Section~\ref{sec:limitations}) its added-member precision is 0.819 against the
rule's 0.518, a paired difference of $+0.301$ [$+0.235$, $+0.366$]; with the day-folded inputs we
first reported, the same comparison gave 0.926 and $+0.408$ [$+0.342$, $+0.479$],
$p_{\mathrm{Holm}}=0.0008$ (Figure~\ref{fig:quiet-precision}, Table~\ref{tab:precision-paired}).
The leak-free comparison was run after the main results and is exploratory.
% source: results/tables/v39_ringfold_precision_paired.csv (proposed_ringfold_vs_b1b, added_member_precision: 0.8190 vs 0.5179, +0.3011 [+0.2349, +0.3657]); v35_precision_paired.csv (proposed_vs_b1b); paper/PREREG_AMENDMENT_2026-11-14.md §3
The matched-ring test cannot see this difference, which is why we test it with the
precision-aware paired bootstrap of Section~\ref{sec:data-stats}. The same test resolves the
gap between the pipeline on Multi-GINe scores and on GFP-XGBoost scores
($+0.139$ ring-F1 [$+0.083$, $+0.196$], $p_{\mathrm{Holm}}=0.0008$), which the matched-ring
test had scored at $p=0.807$: the weaker detector's pipeline matches about as many rings but
adds 2{,}419 edges at 26\,\% precision (day-folded inputs). The policy's ring-F1 gain over the
rule ($+0.027$) does not survive correction ($p_{\mathrm{Holm}}=0.093$); with leak-free inputs
it is $+0.030$ [$+0.010$, $+0.050$], outside the registered family.
% source: results/tables/v35_precision_paired.csv; docs/deep-analysis.md §2; v39_ringfold_precision_paired.csv (proposed_ringfold_vs_b1b, ring_f1: +0.0295 [+0.0101, +0.0501])

\begin{figure}[t]
\centering
\includegraphics[width=0.55\textwidth]{r9_quiet_vs_precision.png}
\caption{Quiet-member recall against added-member precision (day-folded policy inputs, as first
reported; with leak-free inputs the policy's added-member precision is 0.819). The rule and the
policy recover similar shares of quiet members; they differ in how many wrong edges they add.}
% source: results/figures/r9_quiet_vs_precision.png (repository figure)
\label{fig:quiet-precision}
\end{figure}

% Source: results/tables/v35_precision_paired.csv (all eight rows)
\begin{table}[t]
\centering
\caption{Precision-aware paired bootstrap (392 ring-level units, 10{,}000 paired draws, Holm
across four comparisons per metric; primary window, 5 seeds). $p = 0.0002$ raw is the
resolution floor: no draw crossed zero.}
\label{tab:precision-paired}
\small
\begin{tabular}{llccc}
\toprule
Comparison (A vs.\ B) & Metric & $\Delta$ (A $-$ B) [95\,\% CI] & $p_{\mathrm{Holm}}$ & Reject \\
\midrule
proposed vs.\ B1b & added-member prec. & $+0.408$ [$+0.342$, $+0.479$] & $0.0008$ & yes \\ % source: v35_precision_paired.csv (proposed_vs_b1b, added_member_precision)
full vs.\ B1b & added-member prec. & $+0.408$ [$+0.341$, $+0.479$] & $0.0008$ & yes \\ % source: v35_precision_paired.csv (full_vs_b1b, added_member_precision)
GINe full vs.\ GFP-XGB full & added-member prec. & $+0.666$ [$+0.557$, $+0.752$] & $0.0008$ & yes \\ % source: v35_precision_paired.csv (gine_full_vs_gfpxgb_full, added_member_precision)
MEGA full vs.\ GINe full & added-member prec. & $+0.006$ [$-0.042$, $+0.053$] & $0.771$ & no \\ % source: v35_precision_paired.csv (mega_full_vs_gine_full, added_member_precision)
\midrule
proposed vs.\ B1b & ring F1 & $+0.027$ [$+0.001$, $+0.052$] & $0.093$ & no \\ % source: v35_precision_paired.csv (proposed_vs_b1b, ring_f1)
full vs.\ B1b & ring F1 & $+0.034$ [$+0.004$, $+0.065$] & $0.093$ & no \\ % source: v35_precision_paired.csv (full_vs_b1b, ring_f1)
GINe full vs.\ GFP-XGB full & ring F1 & $+0.139$ [$+0.083$, $+0.196$] & $0.0008$ & yes \\ % source: v35_precision_paired.csv (gine_full_vs_gfpxgb_full, ring_f1)
MEGA full vs.\ GINe full & ring F1 & $+0.018$ [$-0.016$, $+0.053$] & $0.303$ & no \\ % source: v35_precision_paired.csv (mega_full_vs_gine_full, ring_f1)
\midrule
\multicolumn{5}{l}{\emph{Leak-free (ring-folded) policy inputs, post hoc; raw $p$, no family correction}}\\
proposed vs.\ B1b & added-member prec. & $+0.301$ [$+0.235$, $+0.366$] & $0.0002^{\P}$ & --- \\ % source: v39_ringfold_precision_paired.csv (proposed_ringfold_vs_b1b, added_member_precision; p_raw 0.0002)
full vs.\ B1b & added-member prec. & $+0.300$ [$+0.233$, $+0.364$] & $0.0002^{\P}$ & --- \\ % source: v45_ringfold_full_precision_paired.csv (full_ringfold_vs_b1b, added_member_precision; p_raw 0.0002)
ring-folded vs.\ day-folded full & added-member prec. & $-0.108$ [$-0.167$, $-0.060$] & $0.0002^{\P}$ & --- \\ % source: v45_ringfold_full_precision_paired.csv (full_ringfold_vs_full)
proposed vs.\ B1b & ring F1 & $+0.030$ [$+0.010$, $+0.050$] & $0.003^{\P}$ & --- \\ % source: v39_ringfold_precision_paired.csv (proposed_ringfold_vs_b1b, ring_f1; p_raw 0.0034)
full vs.\ B1b & ring F1 & $+0.034$ [$+0.008$, $+0.061$] & $0.008^{\P}$ & --- \\ % source: v45_ringfold_full_precision_paired.csv (full_ringfold_vs_b1b, ring_f1; p_raw 0.0082)
ring-folded vs.\ day-folded full & ring F1 & $-0.000$ [$-0.017$, $+0.019$] & $0.96^{\P}$ & --- \\ % source: v45_ringfold_full_precision_paired.csv (full_ringfold_vs_full, ring_f1: -0.0002 [-0.0170, +0.0185], p_raw 0.9617)
\bottomrule
\multicolumn{5}{l}{\footnotesize $^{\P}$ Raw $p$. These draws use 273 units (174 truth rings and 99 off-ring clusters over the methods in that run).}\\
\end{tabular}
\end{table}


\subsection{Community detection is not a stronger baseline than components}
\label{sec:results-baselines}

Replacing connected components with community detection on the same flagged-edge graph does
not help (Table~\ref{tab:ring-methods}, upper block; edge weight is the larger detector score
between two accounts, and the seed fraction is chosen on validation per seed). The best
community method by validation ring-F1, Louvain at resolution 0.5, reaches 0.515 on test,
against 0.521 for seeds and connected components ($-0.006$ [$-0.014$, $+0.000$],
$p_{\mathrm{Holm}}=0.072$) and 0.589 for the full pipeline ($-0.074$ [$-0.118$, $-0.035$],
$p_{\mathrm{Holm}}=0.0004$). Higher resolutions split rings apart (Louvain at resolution 2:
0.107). On the learned policy's expanded subgraph, label propagation in place of components
gives 0.555, below both components (0.582) and the average-linkage partition (0.589). CSGM,
collaborative scatter-gather mining~\citep{tian2025}, is designed for one pattern family and
uses no detector; over five random two-institution splits it reaches 0.083, with ring recall
0.048 over all eight pattern types. These comparisons were added after the main results.
% source: results/tables/v40_external_baselines.json (bar_B: louvain_r0.5 best by val; compare vs b0 / full, paired unit bootstrap 10,000 draws, Holm over two); v45_posthoc_ring_ci.csv (test_visible rows); results/notes/p3_baselines.md

\subsection{Partition: not significant at 174 rings, resolved at 2{,}187}
\label{sec:results-partition}

On HI-Small the partition stage is not significant ($+0.015$ matched rate
[$-0.002$, $+0.036$], $p_{\mathrm{Holm}}=0.49$), nor is partitioning without expansion
($p_{\mathrm{Holm}}=0.23$). On HI-Medium, with 2{,}187 rings, an effect of the same size is
resolved: $+0.015$ [$+0.010$, $+0.020$], $p_{\mathrm{Holm}}=0.0003$
(Section~\ref{sec:results-hm}).
% source: results/tables/v3_closeout_p3.csv (global, full vs b2: mean_delta 0.01486, lo 0.00983, hi 0.01989, p_holm 0.0003)
Its pre-registered hypothesis --- that partitioning would allow a
looser seed budget --- was rejected: validation ring-F1 falls from 0.593 to 0.194 as the seed
fraction grows from 0.10\,\% to 0.50\,\%. A Louvain partition of the same expanded subgraph
leaves ring-F1 unchanged at 0.582.
% source: r9_significance.csv (full_vs_b2, cc_partition_vs_b0); docs/FINDINGS.md (B3, B4); results/tables/r4_partition.csv (b4)
On MEGA-GNN scores, where merged rings are commoner, the gain is the same size ($+0.017$)
and its interval excludes zero ($p_{\mathrm{Holm}} \approx 0.011$): over-merged rings fall
from 9.8 to 8.2 per seed there and stay at 5.2 on ours. The stage removes over-merges only
where they exist, and pays for recall with precision on both detectors.
% source: docs/FINDINGS.md (B3)
Its apparent benefit also depends on two evaluation choices: the recall gain is $+0.024$ at
Jaccard 0.3, $+0.015$ at 0.5 and $+0.001$ at 0.7, and it reverses when only connected pieces
of three or more accounts are scored (Section~\ref{sec:findings-scope}), on both releases.
On HI-Small we do not claim the partition step beyond neutral. Typing, held-out pattern types and time-to-detection are
reported in Appendix~\ref{app:tables}.
% source: docs/FINDINGS.md (B3; dated note, component-level scope)

\subsection{The same pipeline on a larger release}
\label{sec:results-hm}

On HI-Medium the pipeline recovers far less: full-pipeline ring recall is
\ci{0.187}{0.174}{0.200} against HI-Small's 0.461 (Table~\ref{tab:hm-closeout}), 28.7\,\% of rings
are untouched by any seed or expanded edge, and 8.6\,\% are recovered only in part. Its
detector did not meet the adequacy bar fixed before the run, and only two of its three seeds
are usable (Section~\ref{sec:findings-saturation}), so these are not method results on a par
with HI-Small's.
% source: results/tables/v3_closeout_p1.csv (global,full,ring_recall); v3_closeout_verdicts.json (P2_global shares: b_untouched 0.2867, c_under_covered 0.0855); docs/deep-analysis.md §7
Ring-F1 also falls with distance from training: with seeds allocated per day, seed 1 scores
0.408 on test days 6--7 and 0.066 on days 14--15.
% source: results/tables/v3_closeout_blocks.csv (seed 1, per_day, full, ring_f1: block6-7 0.4078, block14-15 0.0656)
What the larger release does provide is power. The partition gain that is not significant on
HI-Small is resolved here at the same size, and the learned policy matches more rings than the
rule ($+0.047$ [$+0.037$, $+0.057$], $p_{\mathrm{Holm}}=0.0003$), where on HI-Small the two were
indistinguishable.
% source: results/tables/v3_closeout_p3.csv (global: full vs b2 +0.0149 [+0.0098, +0.0199]; policy vs b1b +0.0471 [+0.0373, +0.0572]; p_holm 0.0003)

% Source: results/tables/v3_closeout_p1.csv (point / lo / hi), v3_closeout_p3.csv. Intervals: ring bootstrap
% (evaluation/ring_stats.py::bootstrap_over_rings, 1,000 resamples), the same procedure as the HI-Small headline.
\begin{table}[t]
\centering
\caption{Ring recovery on HI-Medium (2{,}187 in-scope rings, test days 6--14, two detector seeds,
in-window members, Jaccard $\ge 0.5$). Brackets: 95\,\% bootstrap intervals over rings.$^{\ddagger}$}
\label{tab:hm-closeout}
\small
\begin{tabular}{llccc}
\toprule
Seed allocation & Method & Ring rec. & Rec. (rings $\ge$3) & Ring F1 \
\midrule
global  & learned expansion     & \ci{0.172}{0.159}{0.185} & \ci{0.243}{0.227}{0.260} & \ci{0.180}{0.166}{0.194} \ % source: v3_closeout_p1.csv (global,proposed)
global  & expansion + partition & \ci{0.187}{0.174}{0.200} & \ci{0.264}{0.246}{0.280} & \ci{0.182}{0.169}{0.195} \ % source: v3_closeout_p1.csv (global,full)
per day & learned expansion     & \ci{0.222}{0.207}{0.238} & \ci{0.315}{0.296}{0.336} & \ci{0.233}{0.217}{0.249} \ % source: v3_closeout_p1.csv (per_day,proposed)
per day & expansion + partition & \ci{0.242}{0.226}{0.257} & \ci{0.340}{0.321}{0.360} & \ci{0.234}{0.219}{0.249} \ % source: v3_closeout_p1.csv (per_day,full)
\bottomrule
\multicolumn{5}{l}{\footnotesize $^{\ddagger}$ Detector seeds 1--2 only (seed 3's scores saturate); the detector did not meet its}\
\multicolumn{5}{l}{\footnotesize pre-registered adequacy bar. Not pooled with HI-Small.}\
\end{tabular}
\end{table}


\subsection{Edge F1 is not a proxy for ring recovery}
\label{sec:results-swap}

We re-ran the ring layer on other detectors at the same fixed 0.1\,\% seed budget, re-choosing
every other downstream parameter on validation per detector and seed (Figure~\ref{fig:edge-ring}, Table~\ref{tab:detector-swap}).
Our MEGA-GNN reproduction reaches $0.6856 \pm 0.0447$ test edge F1, against the
$0.7083 \pm 0.0218$ its authors report for the same configuration; the gap is attributable to
batch-causal sampling (Section~\ref{sec:data}), which their setup does not use\footnote{All MEGA-GNN scores in this paper are produced with a scoring batch of 2048 transactions. The released models were originally scored at 8192, which gives test edge F1 $0.6893 \pm 0.0396$ against $0.6856 \pm 0.0447$ at 2048; per seed, 2048 minus 8192 is $-0.0126$, $-0.0029$, $-0.0013$, $+0.0049$ and $-0.0069$. The causal neighbour sampler cuts batches by size or by 30 minutes of seed time, whichever comes first, so a larger scoring batch lets a transaction see later transactions of the same batch (median look-ahead 180\,s at 2048, 600\,s at 8192). The validation-to-test gap per seed is in Table~\ref{tab:val-test-gap}.}.
% source: results/notes/p7_genagg.md (batch-2048 per-seed table, look-ahead); results/notes/p7_batch_span.json
% Pass 2 (2026-11-16): Table detector_swap and the edge-F1 sentence above now use batch 2048 (r6_detector_swap_sb2048.csv).
% TODO pass 2, still at 8192 and needing recomputation at 2048 before they change: "9.4 points" (2048 gives 0.6856 - 0.5957 = 9.0;
%   also abstract, introduction, conclusion, figure caption); B0 "0.5208 vs 0.5213" (2048: 0.5179); seed-set Jaccard 0.415 and
%   "128.8 and 129.6" rings reached (2048 rings_touched 127.4; seed Jaccard not yet recomputed); fig_edge_vs_ring; scope_sensitivity
%   and component_scope MEGA rows. See results/notes/p5_required_changes.md.
% source: docs/paper/generalisation.md §1

\begin{figure}[t]
\centering
\includegraphics[width=0.62\textwidth]{fig_edge_vs_ring.png}
\caption{Edge F1 against ring F1 for three detectors. From Multi-GINe to MEGA-GNN, 9.4 points
of edge F1 buy no change in ring F1 at the baseline operating point; below Multi-GINe, ring
F1 falls with edge F1.}
% source: results/tables/r6_detector_swap.csv; figure built by paper_draft/make_figs.py
\label{fig:edge-ring}
\end{figure}

% Source: results/tables/r6_detector_swap.csv (GFP-XGBoost, Multi-GINe); results/tables/r6_detector_swap_sb2048.csv (MEGA-GNN rows, scoring batch 2048)
\begin{table}[t]
\centering
\caption{The same ring layer on several detectors at the fixed 0.1\,\% seed budget, every other
downstream parameter re-chosen on validation per detector per seed. With the budget chosen on
validation instead (sensitivity), MEGA-GNN's ring F1 falls to 0.414 (B0) and 0.510 (full); the
Multi-GINe and GFP-XGBoost rows are unchanged. Edge F1 is on the whole test split (mean $\pm$
std over 5 seeds); ring numbers are on the primary window with bootstrap CIs. MEGA-GNN is scored with a batch of 2048 transactions (Section~\ref{sec:results-swap}). Added-member precision here was measured with day-folded expansion-policy inputs; with leak-free (ring-folded) inputs the reference full pipeline's value is 0.817 (Table~\ref{tab:precision-paired}).}
% source (leak-free value): results/tables/v45_ringfold_full_precision_paired.csv (full_ringfold, added_member_precision 0.8174)
\label{tab:detector-swap}
\small
\begin{tabular}{lccccc}
\toprule
Detector & Edge F1 & Ring F1, B0 & Ring F1, full & Quiet rec., full & Added prec., full \\
\midrule
GFP-XGBoost & $0.500 \pm 0.006$ & \ci{0.343}{0.259}{0.426} & \ci{0.451}{0.375}{0.526} & 0.267 & 0.260 \\ % source: r6_detector_swap.csv (gfpxgb rows)
Multi-GINe (shipped) & $0.596 \pm 0.071$ & \ci{0.521}{0.429}{0.607} & \ci{0.589}{0.502}{0.675} & 0.203 & 0.926 \\ % source: r6_detector_swap.csv (gine_causal rows)
MEGA-GNN & $0.686 \pm 0.045$ & \ci{0.518}{0.429}{0.604} & \ci{0.602}{0.517}{0.683} & 0.234 & 0.915 \\ % source: r6_detector_swap_sb2048.csv (mega b0 / full rows, scoring batch 2048)
MEGA-GNN, no embeddings & $0.686 \pm 0.045$ & -- & \ci{0.578}{0.493}{0.659} & 0.200 & 0.787 \\ % source: r6_detector_swap_sb2048.csv (mega_noembed full row)
MEGA-GenAgg & \textit{[TBD]} & \textit{[TBD]} & \textit{[TBD]} & \textit{[TBD]} & \textit{[TBD]} \\ % TODO pass 2: GenAgg placeholder; fill from the five-seed rows once seed 5 is scored
\bottomrule
\end{tabular}
\end{table}


MEGA-GNN outscores Multi-GINe by 9.4 points of edge F1 and produces the same baseline rings
(ring-F1 0.5208 vs.\ 0.5213) at the fixed budget. When validation chooses the budget instead,
two MEGA-GNN seeds pick 0.05\,\% and its rings fall behind (Section~\ref{sec:method}).
% source: results/tables/r6_detector_swap.csv; docs/FINDINGS.md (C6)
This is not the detectors agreeing. Their seed sets overlap at Jaccard 0.415, yet they reach
128.8 and 129.6 of the 174 truth rings. What binds at this operating point is how many rings a
0.1\,\% quantile can reach, not the ranking that fills it. The relation is a ceiling, not
independence: the weaker GFP-XGBoost, 9.6 edge points below, loses 17.7 ring points.
% source: docs/FINDINGS.md (C6)
The claim is conditional on this budget, these detectors and this dataset; a detector
stronger than MEGA-GNN was not tested.

\paragraph{Can the budget be spent better?}
Each GNN detector touches about 74\,\% of truth rings and their touched sets overlap at
71.5\,\%, so their union reaches 150.8 of 174 --- but that is two full budgets. At the same
budget, a best-rank ensemble of the two reaches 133.6, about a fifth of the gap, because 484
of the 864 selected edges are in both detectors' top 0.1\,\%. A per-account budget reaches
fewer (116.2). At an equal budget neither reallocation changes baseline ring-F1 beyond seed
noise (0.521 / 0.515 / 0.519), and under the full pipeline the ensemble's gain
($+0.021 \pm 0.038$) is inside seed spread.
% source: docs/FINDINGS.md (C7, B14 with its dated correction); results/tables/v37_realloc_full.csv

\paragraph{The effective test budget is about half the nominal one.}
Because the quantile is taken over the whole test split, only 463 of 864 seeds land in the
primary window (Section~\ref{sec:findings-tail}). Spending the full budget inside the window
reaches 148.2 rings instead of 129.6 and leaves 28.0 rings untouched instead of 33.2. But
added-member precision (day-folded policy inputs) falls from 0.926 to 0.700, below a 0.85
bound fixed in advance, and
ring-F1 falls from 0.589 to 0.546, on every seed (Table~\ref{tab:window-budget}). The extra
seeds are weaker, and expanding from them costs more rings than the new reach gains. The
budget binds, then, at about 463 in-window seeds, and the reported ring-F1 is not depressed by
the budget lost to the tail. This is a sensitivity analysis on a window examined many times,
and downstream parameters were deliberately not re-tuned.
% source: results/tables/v34_window_budget.json; docs/deep-analysis.md §1

% Source: results/tables/v34_window_budget.json (summary), as tabulated in docs/deep-analysis.md §1
\begin{table}[t]
\centering
\caption{Window-matched seed budget (mean $\pm$ SD over 5 seeds, primary window, 174 truth
rings). Only the seed mask changes; every downstream cell is the shipped, validation-chosen
one. A sensitivity row, not a re-selection: the shipped rule stays the best on ring-F1. Added-member precision here was measured with day-folded expansion-policy inputs; with leak-free (ring-folded) inputs the reference full pipeline's value is 0.817 (Table~\ref{tab:precision-paired}).}
\label{tab:window-budget}
\small
\begin{tabular}{lccc}
\toprule
 & Global (shipped) & In-window quantile & Per-day quantile \\
\midrule
Seeds in window & $463.0 \pm 84.2$ (of 864) & 863 & 864 \\ % source: v34_window_budget.json (seeds_in_window)
Rings touched & $129.6 \pm 4.8$ & $148.2 \pm 6.2$ & $139.8 \pm 7.3$ \\ % source: v34_window_budget.json (rings_touched)
Untouched rings, full & $33.2 \pm 3.6$ & $28.0 \pm 5.5$ & $30.6 \pm 3.8$ \\ % source: v34_window_budget.json (full_b_untouched)
Added-member precision & $0.926 \pm 0.015$ & $0.700 \pm 0.137$ & $0.738 \pm 0.098$ \\ % source: v34_window_budget.json (added_member_precision)
Matched rings, full & $80.2 \pm 5.9$ & $77.6 \pm 5.6$ & $75.8 \pm 5.3$ \\ % source: v34_window_budget.json (full_matched)
Ring F1, full & $0.589 \pm 0.037$ & $0.546 \pm 0.023$ & $0.532 \pm 0.032$ \\ % source: v34_window_budget.json (full_ring_f1)
Ring F1, B0 & $0.521 \pm 0.045$ & $0.503 \pm 0.042$ & $0.485 \pm 0.036$ \\ % source: v34_window_budget.json (b0_ring_f1)
\bottomrule
\end{tabular}
\end{table}


\section{Where ring recovery stops}
\label{sec:results-ceiling}

Figure~\ref{fig:decomposition} accounts for all 174 truth rings. With the full pipeline on the
reference detector, 80.2 are matched per seed, 48.6 have fewer than three visible members and
go unmatched, 33.2 are untouched by any seed or expanded edge, 5.2 are merged with a
neighbour and 6.8 are recovered only in part. Achieved ring recall is
\ci{0.461}{0.393}{0.529}.
% source: docs/FINDINGS.md (B11); results/tables/r6_detector_swap.csv; r9_ring_ci.csv (average, primary, ring_recall)

\begin{figure}[t]
\centering
\includegraphics[width=0.9\textwidth]{fig_ring_decomposition.png}
\caption{What happens to the 174 primary-window truth rings (per seed, full pipeline). A
stronger detector converts untouched rings into matches and into over-merges.}
% source: docs/FINDINGS.md (B11); figure built by paper_draft/make_figs.py
\label{fig:decomposition}
\end{figure}

\paragraph{Small rings are not a hard ceiling.}
60 of the 174 rings have fewer than three visible members, yet the pipeline matches
$11.4 \pm 2.3$ of them per seed, most at Jaccard 1.0, because \texttt{min\_size} applies to
the unrestricted prediction while matching uses the window-restricted one. The share of rings
with three or more visible members (114 of 174) is a reference level, not a bound.
% source: docs/FINDINGS.md (B7, B9, as corrected)

\paragraph{The untouched count belongs to the detector and budget.}
MEGA-GNN at the same budget leaves 24.0 rings untouched instead of 33.2, and 6.4 of the 9.2
newly reached rings become matches. The net matched count rises only from 80.2 to 83.6,
because the stronger detector also merges more rings with their neighbours (5.2 to 8.2).
% source: docs/FINDINGS.md (B8, B11)

\paragraph{Partially recovered rings are clean, short fragments of long patterns.}
In every one of the 34 partially recovered instances the prediction contains no wrong
account, and it holds a median 3 of the ring's 14 members. Thirteen distinct rings produce all
34 instances, half of them STACK.
% source: docs/FINDINGS.md (B12); results/tables/r_under_covered.csv

\paragraph{Most of what is missing is disconnected inside the window.}
For each partially recovered ring we measured the distance from the recovered fragment to the
ring's missing member edges, through the ring itself, inside the primary window. 176 of 216
missing edges (81.5\,\%) lie in a piece with no path to the fragment at all --- all 140 for
STACK, all 24 for BIPARTITE --- and 85.6\,\% were never generated as candidates. The exception
is GATHER-SCATTER, whose 31 missing edges sit one hop away and were declined by the policy; an
oracle admitting exactly those edges completes all 8 instances, about 1.6 rings per seed,
whereas lowering the policy's floor to reach them costs more precision than validation
accepts.
% source: docs/FINDINGS.md (B13, B15); results/tables/r_in_window_connectivity.csv
Widening the window to all ten test days reduces disconnection to 52.3\,\% and removes it for
CYCLE and RANDOM, but STACK stays at 100\,\% (237 of 237): its chains are severed by the
train/validation/test split itself. Traversing the full 18-day graph backwards under the
expansion's own rules reaches none of the 140 missing STACK edges at two or three hops, and a
permutation null reaches none of 216.
% source: docs/FINDINGS.md (B13, B16); results/tables/r_historical_bridge.csv
A supervised-contrastive head~\citep{khosla2020} trained on the frozen detector embeddings to
recognise ring identity does not reconnect the pieces either: same-ring separability across
components rises only from AUC 0.552 to 0.618, against a bar of 0.75 fixed in advance.
% source: docs/deep-analysis.md §10 — registered follow-up; not yet recorded in docs/FINDINGS.md

\paragraph{Observability, counted directly.}
77 of the 174 rings (44\,\%) have no connected in-window piece of three or more accounts. 60
of them are the small rings above; the other 17 have enough visible members, split across
pieces that are each too small, and are scored as ordinary misses. On the 97 rings that do
have a matchable piece, the pipeline recovers 70.6\,\% of the pieces, and STACK recall rises
from 0.036 to 0.493.
% source: docs/FINDINGS.md (dated note, component-level scope); results/tables/v36_component_scope.csv

Taken together, the binding constraint is neither seed allocation nor policy thresholds. It is
how much of a ring is observable and connected inside the evaluation window --- a property of
the split and the scope rule that is easy to mistake for a property of the method.

\subsection{What did not help}
\label{sec:results-negative}
% Pass 2 draft (2026-11-16) from the proposal paragraph in results/notes/m_track.md; descriptive only, no improvement claimed.
We tried five ways of spending the seed budget or reassembling rings after the detector, all at
the fixed 0.1\,\% budget (primary window, test-visible truth, Jaccard 0.5, five seeds). They did not
all use the same reference. Seed fusion and per-account budgets were measured before the leak-free
pipeline existed, against the shipped full pipeline. Account-level seeding was compared with the
shipped full pipeline (0.589). Over-generate-and-verify, motif search and flow-following were
pre-registered against the leak-free full pipeline (0.589 to three decimals; the two pipelines'
ring F1 agree). None raised ring-F1.
\begin{itemize}
\item \emph{Seed fusion and per-account budgets.} A best-rank fusion of two detectors' seeds
gave $+0.021 \pm 0.038$ under the full pipeline, inside seed spread; spreading the budget per
account reached fewer rings and lowered ring-F1.
% source: docs/FINDINGS.md (B14); results/tables/v37_realloc_full.csv
\item \emph{Account-level seeding.} Seeding by account at the same seed count lowered full-pipeline
ring-F1 from \ci{0.589}{0.502}{0.675} to \ci{0.452}{0.377}{0.532} (paired difference $-0.138$
[$-0.184$, $-0.093$]).
% source: results/notes/q_track.md (Q2); results/tables/q_family_holm.csv
\item \emph{Over-generate and verify.} Seeding at twice the budget and filtering with a ring-level
verifier moved ring precision up by $0.098$ [$0.034$, $0.156$] and ring recall down by $0.029$
[$-0.063$, $+0.005$]; the ring-F1 difference, $+0.004$ [$-0.041$, $+0.046$], is
indistinguishable from zero. We report this precision--recall trade as descriptive and secondary.
% source: results/tables/m1_primary.json (other_metrics_paired; primary reading R6 inconclusive, p_Holm 1.00)
\item \emph{Temporal motif search} and \emph{flow-following} from the seeds did not improve on the
seed-only rings. Neither passed its validation gate (best cell 0.569 against 0.585 for the
reference), and their best score-gated cells reproduced the seed-only rings exactly; neither was
run on test.
% source: results/tables/m2_gate.json, m3_gate.json; results/notes/m_track.md (M2, M3)
\end{itemize}
The loss breakdown says why. Of the 174 truth rings in the primary window, 77 have no in-window
piece of three or more accounts (11.4 of them are still matched per seed), 11.0 per seed are
matchable but touched by no seed, 14.4 are recovered too thinly, and merges (1.0) and splits (1.8)
account for little; 68.8 are matched.
% source: results/notes/q_track.md (Q0 table); results/tables/q0_buckets.csv
These counts reconcile with Section~\ref{sec:results-ceiling}: there, the 80.2 matched rings are
$68.8 + 11.4$, and ``untouched'' counts all rings, whereas here it counts only matchable ones.
Reassembly can only rearrange what the seeds already reach, and on this protocol that was not
enough to raise ring-F1.
